\section{The Great Divide and our Ownmost}

In \emph{The Hermetic Tradition}, \textbf{Julius Evola} mentions two competing views of history:

\begin{enumerate}


\item History is the continuous upward evolution of collective humanity.

\item Civilizations arise, mature and die in a series of epochs and disconnected cycles.
\end{enumerate}


The first, he rejects out of hand. The second has some merit yet is inadequate. When making distinctions or categorizing, we need to take care to determine whether the category is real or just an abstraction. The second view, whose premier representative is Oswald Spengler, is an abstraction. That is, it is a mental construct derived from contingent and empirical factors by abstracting out common elements. An abstraction may be useful or convenient for a given purpose, but it is an error when taken as a real category. When this is forgotten, the result will be unclear thinking and the perpetuation of ignorance.

Standing firmly on Rene Guenon's groundbreaking work, \emph{The Crisis of the Modern World}, Evola then makes the distinction between the modern world and the world of Tradition. This distinction is a real metaphysical category as opposed to a mental construct or abstraction. Evola is very precise:

\begin{quotex}
\emph{for the West, we can put the dividing line at the end of the Middle Ages.}
\end{quotex}


For the Traditional man in the West, therefore, what is ``ours'' comprises the Vedic civilization, the ancient Hellenic, Roman and Germanic civilizations and the Medieval civilization. Any attempt to draw the line elsewhere, or to split and oppose those civilizations can be based only on artificial, arbitrary, incomplete and misleading considerations. In particular, the Hermetic Tradition draws on elements from the entirety of those Traditional civilizations; hence it cannot be understood apart from all of them. The failure to embrace the Traditional worldview in its entirety can only lead to confusion. Evola explains:

\begin{quotex}
Who undertakes this study of alchemy and Hermetism without having acquired the ability to rise above the modern mind-set or who has not awakened to a new sensitivity that can place itself in contact with the general spiritual stream that gave life to the tradition in the first place, will succeed only in filing his head with words, symbols, and fantastic allegories.

\emph{The heart of the matter is to re-evoke, by means of an actual transformation of consciousness, the older Traditional basis of understanding and action.}
\end{quotex}


For those of us today, for whom the world of Tradition may still be a dim memory struggling to come to full light, whom do we look to for guidance, who are our ``own''? Besides the Hermetic Tradition and the other figures mentioned in Gornahoor, Evola singles out specifically \textbf{Plotinus} and \textbf{Patanjali} for their metaphysical writings. In his \emph{Saggi sull'Idealismo Magico}, Evola lists as the examples for his philosophical system:

\begin{quotex}
The ``I'', the ``pure act'' would appear as that cosmic centrality that the esoteric indicates, as examples, the type of the \emph{rishi}, the
\emph{yogi}, the \textbf{Christ}, and the \textbf{Buddha}.
\end{quotex}



\flrightit{Posted on 2011-04-19 by Cologero}

\begin{center}* * *\end{center}

\begin{footnotesize}
\begin{sffamily}
\texttt{Mr.Skull on 2011-04-20 at 10:09 said:}

1. What is a `real' category, and how do we reckon it's real? Do we not turn to empirical and contingent factors in reality in order to prove it to be `in reality', in either static or continuous occurrence?

2. Aren't categories, be them real or unreal, ``a mental construct derived from contingent and empirical factors by abstracting out common elements''?

3. I don't see how Evola PROVED Spengler's model to be ``unreal'' or ``an abstraction'', to use his own words (or your interpretation). Disproving Spengler's model doesn't end by claiming it is an abstraction, one's has to show why it's an abstraction or ``unreal'' or ``does not correlate with reality''.

4. Since Evola did not prove Spengler to be wrong, as far as I can tell ``An abstraction may be useful or convenient for a given purpose, but it is an error when taken as a real category.'' can be applied to the claim: ``For the TRADITIONAL man in the West, therefore, what is ``ours'' comprises the Vedic civilization, the ancient Hellenic, Roman and Germanic civilizations and the Medieval civilization.''

I'm no expert on traditionalism, but I've read Evola and Guenon. I enjoyed both writers. Nevertheless, while they do raise some interesting points, they too are entangled within their own verbal `abstractions'.

\hfill

\texttt{Cologero on 2011-04-20 at 13:00 said:}

Nothing in the article said Spengler was ``wrong'', just that it is a theory abstracted from empirical facts. If we take it as absolute, then, because of the discontinuity, then the ancient civilizations can have no relationship to each other, even less so to the modern world. That we reject --- see \textit{The Archeology of the Soul}.

``Real'' categories are based on eternal ideas, not derived from empirical and contingent considerations. We have referred many times to the cosmic order.

Metaphysical concepts are valid always and everywhere, independent of the physical world. To the contrary, the physical world depends on the metaphysical. So they are the polar opposite of a mental abstraction which are products of thought. Metaphysical concepts depend on a higher form of knowing, discussed her many times. Do a search on gnosis, intuition or episteme.

\hfill

\texttt{Graham on 2011-04-20 at 13:35 said:}

Mr. Skull,

\url{http://www.gornahoor.net/?p=1737}


Pay attention to the distinctions between wholes and heaps, and the definitions of the real and the conventional. Tradition is a conventional term used to designate a whole, while Spenglerism is a heap, and nothing but. Wholes cannot be reached by analysis; analysis produces heaps. A heap can be useful, and even conditionally true, but it is still a convention, and fundamentally contingent; the point is to go beyond convention and enter the real. This is the metaphysical `point of view' which Guenon writes of constantly. It's easy for us well-trained critical thinkers to misunderstand these references, and think that he's simply inventing another abstract category, despite how clear and insistent he is on the metaphysical.

\hfill

\texttt{Mr.Skull on 2011-04-20 at 16:26 said:}

Thank you Cologero and Graham for the exegesis.

Cologero,

1. I know the article didn't state Spengler was `wrong'. Nevertheless: A. When Evola critiques, it seems like he's making a perfect abstraction as well. Calling his abstraction a ``metaphysical'' category does not render it real.

B. ``If we take it as absolute, then, because of the discontinuity, then the ancient civilizations can have no relationship to each other, even less so to the modern world.'' Spengler never said anywhere civilizations are disconnected from each other or ``exist in a vacuum''. What Spengler says is that there are civilizations (either Indo-European, or not) that have reached maturity, and these civilizations decline in time according to certain causal laws. He studies the physical relations between human activities, and the decline of civilizations. He never states that these mature civilizations are ``disconnected'' from other, previous epochs or societies, the only thing he may say is that these other epochs and societies / civilizations aren't ``mature'' (he of course defines which civilization qualify to be a mature one). ``Germanic c ivilization'' is not a ``Roman'' C ivilization, but they both may share common elements.

2. Spengler's theory is ``abstract'' as much as Evola's use of metaphysical ideas. Categories derived from abstractions (mental constructs) turn real when they have empirical antecedent, and are initially derived from the physical world, and repeat or cease according to the laws of objective physical reality.

So, ``Real'' categories are based on eternal ideas, not derived from empirical and contingent considerations.'' really means: eternal ideas are based upon empirical and contingent considerations? I can't really see where else can they come from, unless some other, mystical, supernatural force come to exist. Humans induce, than deduce.

Graham: Why does one ought to believe Spengler does not base his theory upon what we can just as well call ``the eternal idea of circularity'', and than demonstrates how the idea is reflected and perpetuated throughout history, in civilizations that reached matureness? Spenglerianism is heavily influenced by the ideas of Plato, and the imminent and WHOLEsome logical consequentiality and causality behind the rise and fall of societies and civilizations. This is derived from the acknowledgment humans are bound to natural laws, which have shaped them as humans and thus their behavior, and that reality works according to this behavioral causality, which gives rise to imminent human constructions (cultural, social, political), ever repeating throughout history, and than declining when X occurs.

\hfill

\texttt{Izak on 2011-04-22 at 02:21 said:}

Skull:

The problem with Spengler's construct is that it was based on an admission of epistemological relativism. I'm quite familiar with Decline of the West. It's a great treatise to say the least. But it has serious problems in logic. For instance, if you look at his second chapter of volume one, on mathematics, he shows a complete rejection of the idea that there is a universal mathematical system rooted in an underlying truth. Instead, he defines mathematical concepts as mere expressions of a cultural worldview. To Spengler, works of art, developments in technology, wars, etc. are only indicators of a culture's stage of development in how it basically feels. Objective calculation becomes lost. While Spengler does present the concept of civilizational rise and decline as a universal rule of law, the great paradox is that he claims that each civilization has a worldview no more valid than the other, and thus that each civilization possesses equal importance to the other. Why more postmodernist thinkers haven't written about this, I haven't the faintest idea. But anyhow, Spengler himself gradually departed from that notion entirely in the interwar years (see: Man and Technics).

The traditionalists differ from Decline in that they refuse to let relativism define their views on civilization. They look toward a series of metaphysical principles as the guiding light, and each civilization to be analyzed either more or less conforms to those principles. There is only one true math and science, and it is a traditional math and science. There is only one true language of symbols, and it comprises the traditional symbols. Etc.

\hfill

\texttt{Mr.Skull on 2011-04-22 at 11:45 said:}

Izak,

What you describe is not necessarily a paradox. Essentially, Spengler says that although ``each culture has its own new possibilities of self-expression'', NO MATTER what view you hold in your civilization, or what set of values you promote and perpetuate which led a civilization to MATURITY, certain laws or ``processes'' which are true and apply to all civilizations, lead to a mature civilizations' exhaustion, old age and decline.

For example, every mature / distinctive civilization eventually stiffen into a type of cultural / social / political ``formalism'' and bureaucracy. Another example: each civilization, once it has finalized the maturity process -- has more residents in dense and large megalopolis, than in the rural areas / countryside.

In the page of concealment of the book ``Man and Techniques'' he writes: ``We are born into this time and must bravely follow the path to the destined end. There is not other way. Our duty is to hold on to the lost position without hope, without rescue... ''. And this is the (pessimistic) essence of Spengler's work.

\hfill

\texttt{Izak on 2011-04-23 at 12:56 said:}

Right, but what I'm saying is that on the one hand, he approaches universal truth as nonexistent. To Spengler, truth is relative to each culture. The traditionalists say, ``There is one truth,'' while Spengler wouldn't say such a thing. You have to remember, he's a Nietzschean.

Man and Technics departs from Decline of the West because Decline talks about how each civilization has just as much importance as the other. It's not ``ethnocentric'' at all; it's completely egalitarian in its view of culture. Man and Technics, however, is Spengler's first attempt to seriously A) talk about mankind's growth and development, and B) analyze the subject of technology. All of the sudden, his relativism is challenged, because he starts to show world history as A) a cohesive whole, and B) completely affected by Western civilization, hence placing more importance on the West.

This is what why the traditionalists aren't as impressed with Spengler. It isn't as though his books are awful, but if you look at his career, you can see changes in his ideas. The traditionalists possess universal principles, and because of that, their principles are much more solid and uncompromised.

\hfill

\texttt{Mr.Skull on 2011-04-23 at 18:00 said:}

``Right, but what I'm saying is that on the one hand, he approaches universal truth as nonexistent.''

And yet, his thesis is completely logical, and it corresponds with both reality and the facts of history, as is the well-known fact ALL humans go through the cycle of ``birth-youth-maturity-old age''. The additional accusations of him being an ``egalitarian'' in his view of culture, evoked by traditionalists has nothing to do with the historiosophic nature of his civilizational analysis, but are completely transparent, since the involvement of the a conception of cultural or value-wise ``importance'' is necessary if the criticism of traditionalists be even slightly ``valid''. By using ``importance'' and the issue of the validity of values, traditionalists can accuse him of invoking ``relativism''. One thing, though, they have to remember: while Spengler never stated any civilization is ``more important'' than the other, he never, explicitly opined each mature civilization is equally ``important'', also -- what is importance? How do we measure it? Important to what? In relation to what?

Spengler also recognized that each of the high civilizations had different values, and thus -- different opportunities and philosophical velocities, WITHIN the confines of existential circularity -- which he examines as a META-ANALYST.

The alleged ``changes'' in his ideas and philosophy, may simply be notional expansions he decided to phrased in order to clarify the ideas expressed in his often misunderstood thesis in DOTW and to give them a historical background. He writes in the preface: ``But experience with the earlier work showed that the majority of readers are not in a position to maintain a general view over the mass of ideas as a whole, and so lose themselves in the detail of this or that domain which is familiar to them.''

In the rest of the book, he uses the same causal method he used in DOTW, to describe how man ``evolves'' from a certain perception of nature and technology, to yet another, finalized one, in which he's possessed by practical technicality and a `mechanical' view of the world.

\hfill

\texttt{Izak on 2011-04-23 at 19:52 said:}

I appreciate your devotion to one of the greatest thinkers of the 20th century, but I can't really agree with much of what you just said.

Man and Technics is a pretty serious departure from Decline of the West. Decline states that there is no such thing as a solid ``human history,'' but that there is only a series of disconnected, autonomous histories of peripheral relevance toward each other. Man and Technics is a look at human history as a cohesive whole, where he begins to lay down very clear principles about human nature, and he examines humanity as an historical concept, not merely a zoological one, which is a clear no-no in Decline.

If you read Prophet of Decline by Farrenkopf, you can see pretty clear and decisive evidence that Spengler was turning around much of his metaphysical ideas during the time leading up to his death. The notes for his book-project on pre-civilized man are very telling.

I think that Spengler's thought was moving into a much better direction in the years up to his death, and it would have been a great contribution to scholarship to see his work come to fruition, but I suppose it was just never meant to be. I also think that Spengler knew quite a lot of things that the traditionalists never have bothered to understand, like economics and political theory... but I think the criticism toward Spengler from their perspective is pretty much right. Spengler never had solid, universal principles beyond his model for civilizational growth and decay, and it caused his philosophy to be useful more as a marker than as a source of metaphysical truth.

Anyhow I'll let you have the last word, since I don't have much more to say on the subject.

\end{sffamily}
\end{footnotesize}

